\section{逐层分解与上同调的有限性}\label{src:A.5}
\begin{epigraph}
MacPherson 告诉我, 我的定理可以看作是如此这般 Grothendieck 如此这般, 这使它变得体面多了.
\epigraphauthor{Jim Propp}
\end{epigraph}
\newterm{逐层分解}{dévissage} 在法语中意为 ``拧开'', 是一种专门用于检验凝聚层性质的 Noether 归纳法. 我们用它把上同调的有限性定理 (\bookref{A.5.3}) 从射影态射推广到固有态射.

\subsection{逐层分解}\label{src:A.5.1-subsection}
\begin{proposition}[逐层分解]\label{src:A.5.1}
设 $X$ 为 Noether 概形, $\mathcal P$ 为 $X$ 上凝聚层的性质, 满足
\begin{enumerate}[(a)]
\item 若 $0\to F'\to F\to F''\to0$ 为 $X$ 上凝聚层的短正合列, 且三者中任意两个满足 $\mathcal P$, 则第三个也满足 $\mathcal P$;
\item 对每个一般点为 $\xi$ 的整闭子概形 $Z\subseteq X$, 存在满足 $\mathcal P$ 的凝聚层 $G$, 使得 $\operatorname{Supp}(G)=Z$, $\mathfrak m_\xi G_\xi=0$, 且 $\dim_{\kappa(\xi)}G_\xi=1$.
\end{enumerate}
则 $X$ 上每个凝聚层都满足 $\mathcal P$.
\end{proposition}
\begin{proof}
对闭子集 $Y\subseteq X$, 若 $X$ 上每个支集满足 $\operatorname{Supp}(F)\subseteq Y$ 的凝聚层 $F$ 都满足 $\mathcal P$, 则称性质 $\mathcal P_Y$ 成立. 用 Noether 归纳法证明命题: 需要证明, 若每个闭子集 $Y'\subsetneq Y$ 都满足 $\mathcal P_{Y'}$, 则 $\mathcal P_Y$ 也成立. 设 $F$ 为 $X$ 上满足 $\operatorname{Supp}(F)\subseteq Y$ 的凝聚层. 给 $Y\subseteq X$ 赋予约化概形结构, 令 $I\subseteq\mathcal O_X$ 为其理想层. 为证明 $F$ 满足 $\mathcal P$, 首先断言, 只需假设 $IF=0$. 由于某个 $n>0$ 满足 $I^nF=0$, 有滤过 $0=I^nF\subseteq I^{n-1}F\subseteq\cdots\subseteq F$ 和短正合列
\[
0\longrightarrow I^{k-1}F/I^kF\longrightarrow F/I^kF\longrightarrow F/I^{k-1}F\longrightarrow0.
\]
由归纳和性质 (a), 只需证明被 $I$ 零化的 $I^{k-1}F/I^kF$ 满足 $\mathcal P$. 其次, 可以假设 $Y$ 不可约, 因而为整概形. 假设 $Y=Y_1\cup Y_2$, $F_1=F|_{Y_1}$, $F_2=F|_{Y_2}$, 则映射 $\phi:F\to F_1\oplus F_2$ 的核与余核支集包含于 $Y_1\cap Y_2$. 对正合列 $0\to\im\phi\to F_1\oplus F_2\to\operatorname{coker}\phi\to0$, 再对 $0\to\ker\phi\to F\to\im\phi\to0$ 应用性质 (a), 可得 $F$ 也满足 $\mathcal P$.

最后, 令 $\xi\in Y$ 为一般点, 性质 (b) 给出满足 $\mathcal P$ 的凝聚层 $G$, 其支集满足 $\operatorname{Supp}(G)\subseteq Y$, 且 $\dim_{\kappa(\xi)}G_\xi=1$. 令 $d=\dim_{\kappa(\xi)}F_\xi$. 由于 $F_\xi$ 和 $G_\xi^{\oplus d}$ 同构, 存在开子概形 $U\subseteq Y$ 和同构 $F|_U\to G^{\oplus d}|_U$. 令 $H$ 为图映射 $F|_U\to F|_U\oplus G^{\oplus d}|_U$ 的像, 并令 $\widetilde H\subseteq F\oplus G^{\oplus d}$ 为 $Y$ 上延伸 $H$ 的子层 \cite[习题 II.$5.15$]{Har77}. 投影 $H\to G^{\oplus d}$ 和 $H\to F$ 在 $U$ 上诱导同构, 因而其核与余核的支集包含于闭子概形 $Y'\subsetneq Y$. 由于 $G$ 满足 $\mathcal P$, 由性质 (a), $G^{\oplus d}$ 也满足 $\mathcal P$. 于是 $H$ 满足 $\mathcal P$, 进而 $F$ 也满足. 另见 \cite[III$_1$.$3.1.2$]{EGA} 和 \cite[Tag 01YI]{SP}.
\end{proof}
\begin{remark}\label{src:A.5.2}
同样的证明还给出一些有用的变体. 首先, 若假设 $\mathcal F^{\oplus n}\in\mathcal P$ 蕴含 $\mathcal F\in\mathcal P$, 则条件 (b) 可以弱化为: 存在满足 $\mathcal P$ 的凝聚层 $G$, 使得 $\operatorname{Supp}(G)\subseteq Z$ 且 $G_\xi\neq0$. 或者, 可把 (b) 换成以下条件: 对每个理想层为 $I_Z$ 的整闭子概形 $Z\subseteq X$, 以及 $X$ 上每个满足 $I_ZF=0$ 的凝聚层 $F$, 存在 $X$ 上满足 $I_ZG=0$ 和 $\mathcal P$ 的凝聚层 $G$, 以及在 $Z$ 的某个非空开子集上为同构的态射 $F\to G$.
\end{remark}

\subsection{逐层分解}\label{src:A.5.2-subsection}
\begin{theorem}[上同调的有限性]\label{src:A.5.3}
设 $f:X\to Y$ 为 Noether 概形之间的固有态射. 对 $X$ 上任意凝聚层 $F$ 和任意 $i\geq0$, $\mathrm R^i f_*F$ 凝聚.
\end{theorem}
\begin{proof}
由平坦基变换 (\bookref{A.2.12}), 可假设 $Y=\Spec A$ 为 Noether 环的谱. 需要证明对 $X$ 上每个凝聚层 $F$ 和任意 $i\geq0$, $\mathrm H^i(X,F)$ 为有限 $A$-模. 当 $X\to\Spec A$ 射影时, 这就是 \cite[定理 III.$5.2$]{Har77}; 其证明很好地展示了上同调的威力, 即使只关心 $\mathrm H^0$ 的情形也是如此. 证明中很快归结为 $X=\mathbb P_A^n$. 选取正合列 $0\to K\to\mathcal O_{\mathbb P_A^n}(-m)^{\oplus d}\to F\to0$, 通过 Čech 上同调计算, 结论对中间项成立; 再由足够高次上同调的消失和降序归纳, 可知结论对 $F$ 成立.

为了应用逐层分解 (\bookref{A.5.1}), 令 $\mathcal P$ 为 $X$ 上凝聚层 $F$ 的以下性质: 对所有 $i$, $\mathrm H^i(X,F)$ 为有限 $A$-模. 它满足 \bookref{A.5.1} 的三取二条件 (a). 为验证 (b), 设 $Z\subseteq X$ 为一般点为 $\xi$ 的整闭子概形. 由 Chow 引理 \cite[习题 II.$4.10$]{Har77}, 存在射影双有理态射 $g:Z'\to Z$, 使得 $Z'\to Z\to Y$ 也射影. 丰沛层 $\mathcal O_{Z'}(1)$ 在 $Z$ 上相对丰沛, 因而对 $d\gg0$ 和 $i>0$, 有 $\mathrm R^i g_*\mathcal O_{Z'}(d)=0$ 和 $\mathrm H^i(Z',\mathcal O_{Z'}(d))=0$. 对 $d\gg0$, 取 $G:=g_*\mathcal O_{Z'}(d)$, 则 $\dim_{\kappa(\xi)}G_\xi=1$. 使用 $\mathrm R^i g_*\mathcal O_{Z'}(d)$ 的消失, Leray 谱序列
\[
\mathrm H^p(Z,\mathrm R^qg_*\mathcal O_{Z'}(d))\Longrightarrow\mathrm H^{p+q}(Z',\mathcal O_{Z'}(d))
\]
意味着 $\mathrm H^i(Z,G)=\mathrm H^i(Z',\mathcal O_{Z'}(d))$ 为有限 $A$-模 (事实上当 $i>0$ 时为 $0$). 另见 \cite[III$_1$.$3.2.1$]{EGA} 和 \cite[Tag 0205]{SP}.

这一论证也可以用导出范畴表述. Chow 引理给出射影双有理态射 $g:X'\to X$, 其中 $X'$ 在 $Y=\Spec A$ 上射影. 考虑正合三角 $F\to\mathrm Rg_*g^*F\to C$. 由于 $g$ 射影, $\mathrm Rg_*g^*F\in D_{\mathrm{Coh}}^b(X)$, 因而 $C\in D_{\mathrm{Coh}}^b(X)$. 由于 $g$ 双有理, $F\to\mathrm Rg_*g^*F$ 在稠密开集上为同构. 使用截断得到的正合三角 $C^n\to C\to\tau_{>n}C$ 以及 $n\ll0$ 时 $C^n=0$, 经归纳可得 $\mathrm R\Gamma(X,C)\in D_{\mathrm{Coh}}^b(Y)$. 由于 $\mathrm R\Gamma(X,g_*g^*F)\in D_{\mathrm{Coh}}^b(Y)$, 得到 $\mathrm R\Gamma(X,F)\in D_{\mathrm{Coh}}^b(Y)$. 将这一论证形式化, 可得到导出范畴的逐层分解版本 (例如 \cite[引理 $15.7$]{LMB00}).
\end{proof}
以下形式函数定理的版本常用于完备局部环 $(A,\mathfrak m)$, 但非局部情形有时也有用.
\begin{theorem}[形式函数]\label{src:A.5.4}
设 $X$ 为 Noether 环 $A$ 上的固有概形, $A$ 关于理想 $I\subseteq A$ 完备. 令 $X_n=X\times_A A/I^{n+1}$. 若 $F$ 为 $X$ 上的凝聚层, 则对每个 $i\geq0$, 有自然同构
\[
\mathrm H^i(X,F)\xrightarrow{\sim}\varprojlim_n\mathrm H^i(X_n,F|_{X_n}).
\]
\end{theorem}
\begin{proof}
见 \cite[定理 III.$11.1$]{Har77} (完备局部环上的射影情形), \cite[定理 $30.8.1$]{Vak17}, \cite[推论 $8.2.4$]{Ill05}, \cite[III$_1$.$4.1.7$]{EGA} 和 \cite[Tag 02OC]{SP}.
\end{proof}
\begin{exercise}\label{src:A.5.5}
更一般地, 证明对 $X$ 上的凝聚层 $F$ 和 $G$, 以及每个 $i\geq0$, 有自然同构
\[
\operatorname{Ext}_{\mathcal O_X}^i(F,G)\xrightarrow{\sim}\varprojlim_n\operatorname{Ext}_{\mathcal O_{X_n}}^i(F|_{X_n},G|_{X_n}).
\]
\end{exercise}
\begin{exercise}\label{src:A.5.6}
使用逐层分解, 把形式函数定理 (\bookref{A.5.4}) 的固有情形归结为射影情形.

\textit{提示:} 归结为证明存在统一的 $d$, 使得对所有 $n,i\geq0$ 有 $\mathrm R^i g_*\mathcal O_{Z'_n}(d)=0$, 其中 $Z'$ 如上同调的有限性 (\bookref{A.5.3}) 的证明中所述, $Z'_n=Z'\times_A A/\mathfrak m^{n+1}$. 为此, 对射影态射 $\Spec_{Z'}\bigoplus_{i\geq0}\mathfrak m^i\mathcal O_{Z'}\to\Spec\bigoplus_{i\geq0}\mathfrak m^i$ 应用 Serre 消失定理 \cite[定理 II.$5.2$]{Har77}.
\end{exercise}
